\section{Confirmation and Chrismation: The Seal of the Holy Spirit}
\label{sec:confirmation}

If Baptism is spiritual birth into the life of Christ, Confirmation---known across the Eastern Churches as Holy Chrismation (\latin{Myron})---is the Pentecostal seal that brings that newborn life to full spiritual stature. Through anointing with episcopally consecrated chrism and the imposition of the hand, accompanied by the words sealing the believer with the Gift of the Holy Spirit, the baptized Christian is bound more closely to the Church, fortified with the sevenfold gifts of the Spirit, stamped with a second indelible character, and empowered as a mature witness to confess and defend the faith of Christ in the world. Both the Latin West and the Christian East confess this single sacrament as the direct bridge linking the rebirth of the font to the sacrificial banquet of the altar.

\subsection{Messianic Anointing, Pentecost, and Apostolic Imposition of Hands}
\label{subsec:confirmation-scripture}

Throughout the Old Testament, anointing with fragrant oil set apart kings, priests, and prophets as consecrated instruments of God's Spirit. The prophets foretold that the coming Messiah---literally \emph{Christos}, the Anointed One---would not merely receive passing impulses of the Spirit, but would carry the Spirit's abiding fullness. Isaias announced this royal and spiritual endowment in the prophecy that supplies the Church's classic enumeration of the seven gifts of the Holy Spirit:
\begin{quotation}
And there shall come forth a rod out of the root of Jesse, and a flower shall rise up out of his root. And the spirit of the Lord shall rest upon him: the spirit of wisdom, and of understanding, the spirit of counsel, and of fortitude, the spirit of knowledge, and of godliness. And he shall be filled with the spirit of the fear of the Lord. (Is~11:1--3)
\end{quotation}
The seventh gift---``the spirit of the fear of the Lord'' (\latin{spiritus timoris Domini})---stands in Is~11:3 alongside the six named in v.~2. Later in the same book, the Servant proclaims: ``The spirit of the Lord is upon me, because the Lord hath anointed me: he hath sent me to preach to the meek, to heal the contrite of heart, and to preach a release to the captives'' (Is~61:1). And the prophet Joel promised that this messianic unction would not remain confined to a single king or prophet, but would overflow upon the whole redeemed people in the final days: ``I will pour out my spirit upon all flesh: and your sons and your daughters shall prophesy'' (Joel~2:28; cf.~Ezech~36:27).

In the fullness of time, this prophecy came to pass upon Jesus of Nazareth. Conceived by the Holy Spirit in Mary's womb, He was visibly anointed at the Jordan when the Spirit descended and remained upon Him as the Father named Him Beloved Son (Matt~3:16--17). Standing in the synagogue at Nazareth, Jesus read Is~61:1 aloud, rolled up the scroll, and declared: ``This day is fulfilled this scripture in your ears'' (Luke~4:18--21). Yet the Spirit resting upon the Head was destined to be poured forth upon the entire Body once the Paschal Mystery reached its consummation. At the Feast of Tabernacles, Jesus cried out: ``If any man thirst, let him come to me and drink. He that believeth in me, as the scripture saith: Out of his belly shall flow rivers of living water.'' St.~John immediately explains: ``Now this he said of the Spirit which they should receive who believed in him: for as yet the Spirit was not given, because Jesus was not yet glorified'' (John~7:37--39). At the Last Supper, Christ promised ``another Paraclete'' who would abide with the disciples forever (John~14:16--17), and before ascending into heaven He commanded them to await the promised power: ``you shall receive the power of the Holy Ghost coming upon you, and you shall be witnesses unto me in Jerusalem, and in all Judea, and Samaria, and even to the uttermost part of the earth'' (Acts~1:8).

On the morning of Pentecost, rushing wind and parted tongues of fire filled the upper room, transforming fearful disciples behind locked doors into fearless heralds of the risen Lord (Acts~2:1--4). Crucially, the Acts of the Apostles reveals that this Pentecostal outpouring was not a one-time marvel restricted to the Twelve, but an enduring gift imparted through a distinct apostolic rite following water Baptism. When Philip the deacon preached in Samaria and baptized the new believers, the Apostles in Jerusalem sent Peter and John to complete their initiation:
\begin{quotation}
Now, when the apostles, who were in Jerusalem, had heard that Samaria had received the word of God, they sent unto them Peter and John. Who, when they were come, prayed for them that they might receive the Holy Ghost. For he was not as yet come upon any of them: but they were only baptized in the name of the Lord Jesus. Then they laid their hands upon them: and they received the Holy Ghost. (Acts~8:14--17)
\end{quotation}
The same two-step rhythm recurs at Ephesus, where St.~Paul encounters twelve disciples who knew only John's baptism of repentance. After instructing them in Jesus, ``they were baptized in the name of the Lord Jesus. And when Paul had imposed his hands on them, the Holy Ghost came upon them: and they spoke with tongues and prophesied'' (Acts~19:5--6).

In the epistles, the apostolic Churches already speak of this rite using the permanent vocabulary of strengthening, anointing, and sealing. St.~Paul writes to the Corinthians: ``Now he that confirmeth us with you in Christ and that hath anointed us, is God: Who also hath sealed us and given the pledge of the Spirit in our hearts'' (\latin{Qui autem confirmat nos vobiscum in Christo, et qui unxit nos Deus: qui et signavit nos, et dedit pignus Spiritus in cordibus nostris}, 2~Cor~1:21--22). To the Ephesians he recalls that after believing the Gospel, ``you were signed with the holy Spirit of promise, Who is the pledge of our inheritance'' (Eph~1:13--14). And the Letter to the Hebrews lists ``the doctrine of baptisms and imposition of hands'' side by side among the foundational elements of Christian initiation (Heb~6:1--2). Taken together, Acts~8, Acts~19, 2~Cor~1:21--22, Eph~1:13--14, and Heb~6:2 display the unmistakable apostolic roots of Confirmation: a sacramental communication of the Holy Spirit distinct from Baptism, ordered to apostolic strength, and signified by the laying on of hands and spiritual anointing.

\subsection{The Fathers on the Royal Unction and Spiritual Seal}
\label{subsec:confirmation-fathers}

In the early Church, Baptism and the post-baptismal anointing with imposition of hands were ordinarily celebrated in one continuous Paschal liturgy, yet the Fathers spoke clearly of the distinct spiritual gift bestowed in the chrismation that followed rising from the font. Two fourth-century masters---one from Jerusalem in the East, one from Milan in the West---allow us to hear the Church's mystagogical preaching directly.

In the third \emph{Mystagogical Catechesis}, titled \emph{On Chrism} (traditionally ascribed to St.~Cyril of Jerusalem, with modern scholarship also considering his successor John~II of Jerusalem), the preacher turns immediately from the baptismal pool to the sacred ointment applied to the newly illuminated. Just as Christ received the descent of the Spirit after emerging from the Jordan, the baptized receive their share in His messianic unction:
\begin{quotation}
And to you in like manner, after you had come up from the pool of the sacred streams, there was given an Unction, the anti-type of that wherewith Christ was anointed; and this is the Holy Ghost. \dots\ But beware of supposing this to be plain ointment. For as the Bread of the Eucharist, after the invocation of the Holy Ghost, is mere bread no longer, but the Body of Christ, so also this holy ointment is no more simple ointment, nor (so to say) common, after invocation, but it is Christ's gift of grace, and, by the advent of the Holy Ghost, is made fit to impart His Divine Nature. Which ointment is symbolically applied to thy forehead and thy other senses; and while thy body is anointed with the visible ointment, thy soul is sanctified by the Holy and life-giving Spirit. (\emph{Mystagogical Catechesis}~III.1, 3)
\end{quotation}
The preacher then walks the neophytes through each anointed sense---the forehead to erase the shame of the first transgressor, the ears to hear divine mysteries, the nostrils to breathe the sweet fragrance of Christ, the breast as a breastplate of righteousness against the devil's wiles---and lands on an insight of startling simplicity: ``Having been counted worthy of this Holy Chrism, ye are called Christians, verifying the name also by your new birth'' (\emph{Mystagogical Catechesis}~III.4--5). Because \emph{Christos} means the Anointed One, it is in the sacred Chrism that the Christian visibly receives the royal, priestly, and prophetic unction that justifies his very name.

In the Latin West, St.~Ambrose of Milan places the same mystery before his neophytes in \emph{De mysteriis}~VII.42, joining the prophetic sevenfold Spirit of Is~11:2--3 directly to the Pauline seal of 2~Cor~1:21--22:
\begin{quotation}
And then remember that you received the seal of the Spirit; the spirit of wisdom and understanding, the spirit of counsel and strength, the spirit of knowledge and godliness, and the spirit of holy fear \dots\ God the Father sealed you, Christ the Lord strengthened you, and gave the earnest of the Spirit in your heart, as you have learned in the lesson from the Apostle. (St.~Ambrose, \emph{De mysteriis}~VII.42)
\end{quotation}
Here the two great biblical signs converge without rivalry. The Latin tradition gradually highlighted the apostolic laying on of hands (\latin{manus impositio}) recalled in Acts~8 and~19 alongside the episcopal signing of the forehead with chrism, naming the sacrament \emph{Confirmation} (\latin{confirmatio}) for its strengthening grace; the Greek and Eastern Syriac traditions highlighted the royal and Pentecostal anointing with sacred \emph{myron} across the principal senses, naming the sacrament \emph{Chrismation} (\latin{myron}, \latin{sphragis}, the Seal of the Gift of the Holy Spirit). Both East and West confessed the same divine reality: the personal gift of the Holy Spirit sealing the baptized soul in the likeness of the Anointed Christ.

\subsection{Sacramental Constitution: Matter, Form, Minister, and Subject}
\label{subsec:confirmation-constitution}

The Council of Florence (\emph{Exsultate Deo}, 1439) and the Council of Trent (Session~VII, Canons on Confirmation) solemnly defined that Confirmation is a true and proper sacrament instituted by Christ---not a mere catechetical ceremony or renewal of baptismal vows---and defended the spiritual efficacy of consecrated chrism and the ordinary ministry of bishops. St.~Thomas Aquinas (\emph{ST}~III, q.~72), the \emph{Roman Catechism} (Part~II), the Second Vatican Council (\emph{Lumen Gentium}~11; \emph{Orientalium Ecclesiarum}~12--14), Pope Paul~VI's Apostolic Constitution \emph{Divinae Consortium Naturae} (15~August 1971), and the 1983 and 1990 Codes of Canon Law set forth its exact sacramental constitution across the Catholic communion.

\paragraph{Matter.}
The remote matter of Confirmation is \textbf{sacred chrism} (\latin{chrisma} / \latin{myron}): olive oil---or, where lawful under modern liturgical norms, another plant oil---blended with aromatic balsam or fragrant spices and consecrated beforehand by a bishop in the Latin Church (CIC~1983, c.~880~\S2), or by a Patriarch, Major Archbishop, or bishop in the Eastern Catholic Churches (CCEO~1990, c.~693). St.~Thomas Aquinas explains the deep fitness of this two-part matter (\emph{ST}~III, q.~72, a.~2). Olive oil signifies the radiant, supple strength of the Holy Spirit---the oil with which athletes anointed their limbs before combat and lamps burned with clear light---while sweet-smelling balsam signifies both the ``good odour of Christ'' (2~Cor~2:15) diffused by a mature Christian's public witness and the incorruptibility by which grace preserves the soul from spiritual decay.

The proximate matter in the Latin Church, definitively settled by Pope Paul~VI in \emph{Divinae Consortium Naturae} (1971) and codified in CIC~1983, c.~880~\S1, is the \textbf{anointing with chrism on the forehead done by the laying on of the hand}---that is, the physical imposition of the minister's hand integrated into the very act of signing the forehead with the thumb dipped in chrism (following the solemn extended-hands prayer invoking the sevenfold Spirit over the candidates). This clarifies centuries of scholastic discussion about how imposition of hands and chrismation relate: both biblical gestures are physically united in the essential act. In the Byzantine and other Eastern Churches, the proximate matter is the anointing with consecrated \emph{myron} in the sign of the cross upon the forehead and the other principal senses (eyes, nostrils, mouth, ears, breast, hands, and feet in Byzantine usage), sealing the whole bodily person for the Spirit's service (CCEO~1990, c.~692; CCC~1299--1300).

\paragraph{Form.}
In \emph{Divinae Consortium Naturae} (1971), Pope Paul~VI renewed the Latin essential form by adopting an ancient formula rooted in the Byzantine tradition, expressing the Pentecostal Gift with crystal simplicity (cf.~CCC~1300):
\begin{itemize}
  \item \textbf{Latin form:} ``N., be sealed with the Gift of the Holy Spirit'' (\latin{Accipe signaculum doni Spiritus Sancti}).
  \item \textbf{Byzantine Eastern form:} ``The seal of the gift of the Holy Spirit'' (\latin{Sphragis doreas Pneumatos Hagiou}; CCC~1300---rendered ``the seal of the gift that is the Holy Spirit'' in the Vatican web English edition, whereas the USCCB second edition renders ``the seal of the gift of the Holy Spirit''; other Eastern rites employ closely cognate approved formulas invoking the seal of the Spirit).
\end{itemize}
The older medieval Latin form expounded by St.~Thomas (\emph{ST}~III, q.~72, a.~4) and the \emph{Roman Catechism}---``I sign thee with the sign of the cross, and I confirm thee with the chrism of salvation, in the name of the Father, and of the Son, and of the Holy Ghost''---signified the same grace of fortification and cross-bearing witness. By restoring the concise sealing formula, the Church highlighted the shared apostolic faith uniting East and West: what is given is not merely an abstract spiritual boost, but the Personal Gift (\latin{Donum}) who is the Holy Spirit Himself.

\paragraph{Minister and Intention.}
No sacrament exhibits the complementarity of Latin and Eastern Catholic discipline more vividly than the ministry of Confirmation and Chrismation.

In the \textbf{Latin Church}, the \textbf{bishop} is the ordinary minister of Confirmation (Trent, Sess.~VII, can.~3 on Confirmation; CIC~1983, c.~882). St.~Thomas grounds this reservation directly in Acts~8:14--17: just as in the primitive Church the fullness of the Holy Spirit was imparted through the Apostles Peter and John after Philip the deacon had baptized, so the conferring of Confirmation belongs ordinarily to the bishops who succeed the Apostles and hold supreme pastoral authority in their local Churches (\emph{ST}~III, q.~72, a.~11). Seeing the bishop impose hands and sign the forehead visibly links the newly confirmed Christian to the apostolic college and the unity of the diocesan Church. At the same time, the Church provides generously for presbyteral administration where pastoral need or the unity of initiation calls for it. Under CIC~1983, c.~883, a presbyter validly administers Confirmation by universal law itself when he is equivalent in law to a diocesan bishop, when he baptizes a person beyond infancy or receives an already baptized person into full Catholic communion, or for any baptized person in danger of death (where every priest possesses the faculty immediately). Beyond those cases, the diocesan bishop may delegate or associate presbyters by special faculty when necessity requires (CIC~1983, c.~884).

In the \textbf{Eastern Catholic Churches}, ancient discipline preserved unbroken from the patristic age makes the \textbf{presbyter} an ordinary and regular minister of Chrismation, ordinarily conferring it immediately after Baptism in a single celebration. The Second Vatican Council solemnly confirmed and restored this ancestral Eastern patrimony in \emph{Orientalium Ecclesiarum}~13--14, affirming that Eastern presbyters validly administer Chrismation with chrism blessed by a Patriarch or bishop, not only to Eastern faithful but even to Latin faithful under licit canonical norms. The \emph{Code of Canons of the Eastern Churches} codifies this discipline in cc.~693--696: the presbyter baptizes and chrismates conjointly, while the \textbf{consecration of the sacred \emph{myron} remains reserved to the Patriarch or bishop} (CCEO, c.~693). Thus Eastern Chrismation preserves the exact same apostolic and episcopal bond as Latin Confirmation, expressing it through the episcopally consecrated oil carried from the Patriarchal or eparchial cathedral to every parish font, rather than by postponing the rite until a bishop can visit in person. In every tradition, the minister must act with at least the intention to do what the Church does.

\paragraph{Capable Subject.}
Every living, validly baptized person who has not yet been confirmed is capable of receiving Confirmation (CIC~1983, c.~889~\S1; CCEO~1990, c.~692). Because Baptism is the indispensable foundation of all supernatural life, an unbaptized person cannot be validly confirmed. For fruitful reception, a candidate who has reached the use of reason must be in the state of sanctifying grace, suitably instructed, and able to renew his baptismal promises (CIC~1983, c.~889~\S2); if conscious of mortal sin, he must first receive sacramental Confession. In danger of death, even an uninstructed infant or young child in the Latin Church may and should be confirmed (CIC~1983, cc.~883 n.~3, 891), for divine grace does not wait upon human eloquence.

\begin{fourstable}{Constitutive Dimension}{Common Catholic Faith}{Latin Church (CIC~1983)}{Eastern Churches (CCEO~1990)}
Matter & Anointing with episcopally consecrated chrism / \latin{myron} with imposition of hand & Chrism blessed by bishop; anointing of forehead with laying on of hand (c.~880; \emph{Divinae Consortium Naturae}) & \latin{Myron} blessed by Patriarch or bishop (c.~693); anointing of forehead and principal senses (c.~692)\\
Form & Sealing words naming the Gift of the Holy Spirit (cf.~CCC~1300) & ``N., be sealed with the Gift of the Holy Spirit'' & Byzantine: ``The seal of the gift of the Holy Spirit'' (cognate approved Eastern formulas)\\
Minister & Apostolic episcopal bond preserved in both rites & Ordinary: bishop (c.~882); presbyter by law for adult initiation / danger of death (c.~883) or delegation (c.~884) & Bishop or presbyter regularly with Patriarchal / episcopal \latin{myron} (\emph{OE}~13--14; cc.~694--696)\\
Subject and Timing & Validly baptized, unconfirmed living person; never repeated & Age of discretion or conference age (c.~891); adults immediately after Baptism (c.~866); danger of death at any age & Conjointly with Baptism in infancy and adulthood outside true necessity (c.~695~\S1)\\
\end{fourstable}

\subsection{Character, Sevenfold Grace, Change, and Order of Ends}
\label{subsec:confirmation-effects}

Confirmation unfolds the same threefold sacramental structure as Baptism (\emph{ST}~III, q.~72, aa.~5, 7):
\begin{itemize}
  \item \textbf{\latin{Sacramentum tantum}} (the outward sign alone): the anointing with consecrated chrism joined to the imposition of the hand and the sealing words.
  \item \textbf{\latin{Res et sacramentum}} (the intermediate reality and sign): the indelible \emph{Confirmation character} stamped upon the baptized soul.
  \item \textbf{\latin{Res tantum}} (the ultimate proper grace): a special Pentecostal outpouring of the Holy Spirit, an increase and deepening of sanctifying grace, fuller vigor of the sevenfold gifts of the Spirit, and supernatural fortitude for public Christian witness.
\end{itemize}

\paragraph{The Confirmation Character.}
Like Baptism and Holy Orders, Confirmation imprints an indelible spiritual character upon the soul and therefore can never be repeated (Trent, Sess.~VII, can.~9 on Sacraments in General; \emph{ST}~III, q.~72, a.~5; CIC~1983, c.~845~\S1; CCEO~1990, c.~672~\S1). St.~Thomas asks why a second character is needed when Baptism already configures the soul to Christ's priesthood. He answers from the analogy of human growth: in Baptism a man receives spiritual life as a newborn child, receiving the passive power to undergo Christ's saving actions for his own sanctification; in Confirmation he arrives at ``the perfect age, as it were, of the spiritual life'' (\emph{ST}~III, q.~72, a.~1), quoting the ancient saying attributed to Melchiades: ``In Baptism we are born again unto life; after Baptism we are strengthened.'' The Confirmation character is a distinct spiritual potency empowering the baptized Christian ``to do those things which pertain to the spiritual combat with the enemies of the Faith'' (\emph{ST}~III, q.~72, a.~5)---to confess the name of Christ publicly, openly, and without fear, as a deputized witness of the Gospel.

\paragraph{The Pentecostal Outpouring and Sevenfold Gifts.}
In an unobstructed recipient, Confirmation brings a genuine increase of sanctifying grace together with a new invisible mission of the Holy Spirit for strength (\emph{ST}~III, q.~72, a.~7). Because the baptized soul already possesses sanctifying grace and the seven gifts in seed from the font, Confirmation does not introduce a different species of divine life; rather, it expands, deepens, and activates those gifts toward public apostolic action. The Second Vatican Council summarizes this ecclesial commissioning in \emph{Lumen Gentium}~11: bound more perfectly to the Church through Confirmation, the faithful are enriched with a special strength of the Holy Spirit and are thereby more strictly obliged to spread and defend the faith by word and deed as true witnesses of Christ.

\begin{threestable}{Gift of the Spirit (Is~11:2--3)}{Supernatural Perfection of the Soul}{Fruit in Confirmed Christian Witness}
Wisdom (\latin{sapientia}) & Judges all created things from the highest divine causes through loving connaturality with God & Enables the believer to weigh worldly pressures and honors against the truth of the Cross\\
Understanding (\latin{intellectus}) & Penetrates deeply into revealed mysteries beyond surface words & Illuminates Scripture, creed, and liturgy so the Christian confesses the faith with clarity\\
Counsel (\latin{consilium}) & Directs practical judgment swiftly and rightly in complex concrete choices & Guides courageous fidelity when conscience faces social, familial, or civic conflict\\
Fortitude (\latin{fortitudo}) & Sustains the will steadfastly through fear, ridicule, hardship, and even martyrdom & Supplies the defining Pentecostal boldness (\latin{parrhesia}) to confess Christ without shame\\
Knowledge (\latin{scientia}) & Sees created realities accurately in their true order toward God & Guards the Christian from both worldliness and false spiritualism in daily vocations\\
Piety / Godliness (\latin{pietas}) & Moves the heart with filial reverence toward God the Father and fraternal mercy toward all & Turns defense of the faith away from harsh polemic into loving apostolic service\\
Fear of the Lord (\latin{timor Domini}) & Inspires reverent awe of God's majesty and horror of offending His love by sin (Is~11:3) & Anchors public courage in humility, keeping the witness faithful to truth rather than vanity\\
\end{threestable}

\paragraph{The Metaphysical Change Effected.}
Metaphysically, Confirmation produces neither a substantial change (which belongs to the Eucharist alone) nor a second spiritual birth, since the subject is already reborn in Christ. It effects a real \textbf{accidental formal strengthening and elevation} of the baptized soul---imprinting the second indelible character as a permanent spiritual potency in the intellective faculty and increasing habitual sanctifying grace with its proper sacramental assistance for steadfast witness---paired with a deepened \textbf{ecclesial relation} joining the believer more closely to the Church's apostolic mission. If an adult approaches Confirmation validly while knowingly in unrepented mortal sin (\latin{obex}), the Confirmation character is truly impressed, while the fruitful grace is blocked until restored through Penance (\emph{reviviscentia}).

\paragraph{The Order of Ends.}
The ends of Confirmation stand in clear hierarchical order:
\begin{enumerate}
  \item \textbf{Primary proper end:} spiritual maturity in Christ and fortification by the special outpouring of the Holy Spirit for public confession of the faith and spiritual combat.
  \item \textbf{Intrinsic secondary ends:} closer union with the Church and the apostolic episcopate, deepening of baptismal adoptive filiation, and fuller readiness for conscious, active participation in the Eucharistic sacrifice.
  \item \textbf{Contingent ends:} in Latin celebrations at the age of discretion or adolescence, the personal renewal of baptismal promises, catechetical consolidation, and visible strengthening of parish communion during the bishop's pastoral visitation.
  \item \textbf{Ultimate end:} eternal glory and perfect union with the Father and the Son in the love of the Holy Spirit, of which the Spirit's seal is already the living pledge (2~Cor~1:22; Eph~1:14).
\end{enumerate}

\subsection{The Unity of Initiation Across East and West}
\label{subsec:confirmation-unity}

A faithful Catholic account of Confirmation must honor how the twenty-four Catholic Churches \latin{sui iuris} arrange its timing within the single triad of Christian Initiation. In the first Christian centuries, Baptism, Chrismation, and First Eucharist formed one seamless Paschal night for adults and infants alike. As parishes multiplied across vast rural dioceses in the Latin West during the early Middle Ages, bishops could no longer be present at every infant baptism. Wishing to preserve the bishop's personal ministry of Confirmation, the Latin Church permitted parish priests to baptize infants promptly while waiting for the bishop's visitation to confer Confirmation. Over time, this pastoral arrangement produced a separated sequence for infants baptized in the Latin Church---Baptism in infancy, First Eucharist upon reaching the age of discretion (CIC~1983, cc.~913--914), and Confirmation either around the age of discretion or at a later age set by the Episcopal Conference (CIC~1983, c.~891).

The Christian East took the complementary path to preserve the very same apostolic bond. Keeping Baptism, Chrismation, and Eucharist united in a single liturgy even for infants, the Eastern Churches entrusted Chrismation to the baptizing presbyter while reserving the consecration of the sacred \emph{myron} to the Patriarch or bishop (\emph{Orientalium Ecclesiarum}~13--14; CCEO~1990, cc.~693--697, 710). Where historical Latinization led several Eastern Catholic Churches to delay First Eucharist or separate Chrismation, the Second Vatican Council, the 1990 Eastern Code, and the Congregation for the Eastern Churches' 1996 \emph{Instruction for Applying the Liturgical Prescriptions of the CCEO} (\S\S42--51) encouraged Synods and Councils of Hierarchs to recover their ancestral organic unity of initiation.

Meanwhile, within the Latin Church itself, adult initiation strictly maintains the unified order of Baptism, Confirmation, and Eucharist in one celebration at the Easter Vigil (CIC~1983, c.~866), and a number of dioceses exercise the norm of CIC~1983, c.~891 by restoring Confirmation prior to First Holy Communion at the age of discretion. Each lawful discipline highlights a true dimension of the sacrament: the unified Baptism--Confirmation--Eucharist sequence manifests the indivisible theological flow from the font through the Spirit's seal to the altar, while the Latin episcopal celebration at a later age visibly underscores communion with the bishop and the conscious personal commitment of a growing believer. Yet Confirmation is never a ``graduation'' from religious education or a merely human coming-of-age ceremony; whether conferred upon an infant wrapped in white linen or upon a youth or adult kneeling before the bishop, its essence is entirely divine---the sovereign Gift of the Holy Spirit sealing the baptized Christian for the Eucharist and for the mission of the Gospel.
